\documentclass[aspectratio=169]{beamer}
\usepackage[spanish]{babel}
\usepackage[utf8]{inputenc}
\usepackage{amsmath,amssymb}
\usepackage{graphicx}
\usepackage{xcolor}
\usepackage{booktabs}
\usepackage{enumitem}

% ===== Colores institucionales UdeC =====
\definecolor{UdeCblue}{RGB}{0,51,153}
\definecolor{UdeCgold}{RGB}{255,204,0}
\setbeamercolor{structure}{fg=UdeCblue}
\setbeamercolor{frametitle}{bg=UdeCblue!10, fg=UdeCblue}

\usetheme{Madrid}
\usecolortheme{default}
\setbeamertemplate{navigation symbols}{}
\setbeamertemplate{footline}[frame number]
\setbeamertemplate{itemize items}[circle]

% ===== Título =====
\title[\textbf{Ayudantía 6 – IS–LM y DA}]{\textbf{Ayudantía 6 – Modelo IS–LM y Demanda Agregada}}
\author{\textbf{Profesor:} Cristian Mardones, Ph.D.\\[2pt]\textbf{Ayudante:} Ignacio Garrido Urra}
\institute{Facultad de Ingeniería\\Departamento de Ingeniería Industrial – Universidad de Concepción}
\date{Semestre 2025–2}

% ===== Comando para espacio de desarrollo =====
\newcommand{\espaciodesarrollo}{
\vspace{0.6cm}
\begin{beamercolorbox}[wd=\textwidth,rounded=true,shadow=false,sep=8pt,center]{block body alerted}
\vspace{1.8cm}
\textit{(Espacio para desarrollo del ejercicio)}
\vspace{1.8cm}
\end{beamercolorbox}
}

\begin{document}

% ===== Portada =====
\begin{frame}
  \titlepage
\end{frame}



\begin{frame}{Repaso Teórico: Curva IS}
\begin{block}{Equilibrio del mercado de bienes}
\[
Y = C(Y - T) + I(i) + G
\]
\[
C = C_0 + c(Y - T) \quad \Rightarrow \quad i = \frac{1}{b}\,[A - (1 - c(1-t))Y]
\]
\end{block}
\begin{itemize}
  \item Pendiente negativa: \(i \downarrow \Rightarrow I \uparrow \Rightarrow Y \uparrow\).
  \item Desplazamientos:
    \begin{itemize}
      \item \(G \uparrow\) o \(T \downarrow \Rightarrow\) IS se desplaza a la derecha.
      \item \(c \uparrow\) aumenta la sensibilidad de \(Y\) a cambios en gasto autónomo.
    \end{itemize}
\end{itemize}
\end{frame}

\begin{frame}{Repaso Teórico: Curva LM}
\begin{block}{Equilibrio del mercado monetario}
\[
\frac{M}{P} = L(Y,i) = kY - hi
\Rightarrow i = \frac{k}{h}Y - \frac{1}{h}\frac{M}{P}
\]
\end{block}
\begin{itemize}
  \item Pendiente positiva: \(Y \uparrow \Rightarrow i \uparrow\).
  \item Si \(M/P \uparrow \Rightarrow\) LM se desplaza hacia la derecha (menores tasas).
  \item Si \(P \uparrow \Rightarrow (M/P) \downarrow \Rightarrow\) LM se desplaza a la izquierda.
\end{itemize}
\end{frame}

\begin{frame}{De IS–LM a Demanda Agregada (DA)}
\begin{itemize}
  \item A cada nivel de precios \(P\), se determina el equilibrio IS–LM.
  \item Al aumentar \(P\), la oferta monetaria real \((M/P)\) cae → LM se desplaza a la izquierda → \(Y\) cae.
  \item Resultado: relación inversa entre \(P\) y \(Y\) → pendiente negativa de la curva DA.
\end{itemize}
\end{frame}

% ===== Ejercicios =====

\begin{frame}{Ejercicio 1: Curvas IS, LM y DA}
Obtenga las expresiones matemáticas de la curva IS y de la curva LM, y derive la curva DA.
\espaciodesarrollo
\end{frame}

\begin{frame}{Ejercicio 1}
    
\end{frame}

\begin{frame}{Ejercicio 2: Problema 3 – Certamen 2013}
\[
C = 60 + 0.8Y_D, \quad t = 0.2, \quad I = 70, \quad G = 230, \quad TR = 110
\]
\begin{itemize}
  \item Calcule el nivel de ingreso de equilibrio.
  \item Determine el multiplicador del gasto.
  \item Calcule el balance fiscal.
  \item Analice el cambio si \(t = 0.3\).
\end{itemize}
\espaciodesarrollo
\end{frame}
\begin{frame}{Ejercicio 2}
    
\end{frame}


\begin{frame}{Ejercicio 3: Problema 3 – Certamen 2014}
Discuta la siguiente afirmación:
\begin{quote}
\textit{El alza del impuesto a la renta disminuiría el PIB (\(\Delta Y / \Delta t < 0\)).}
\end{quote}
\textbf{Considere:} efecto sobre \(Y_D\), consumo, pendiente y desplazamiento de IS.
\espaciodesarrollo
\end{frame}

\begin{frame}{Ejercicio 3}
    
\end{frame}

\begin{frame}{Ejercicio 4: Equilibrio IS–LM}
\[
C = 140 + 0.3Y_D,\quad I = 1000 - 100i,\quad G = 800,\quad TR = 250
\]
\[
t = 0.2,\quad L = 0.25Y - 250i,\quad \tfrac{M}{P}=300
\]
\begin{enumerate}[label=\alph*)]
  \item Ecuación de IS y LM.  
  \item Tasa de interés e \(Y\) de equilibrio.  
  \item Impacto de \(G \uparrow\).  
  \item Impacto de \(P \uparrow\).  
\end{enumerate}
\espaciodesarrollo
\end{frame}
\begin{frame}{Ejercicio 4}
    
\end{frame}

\begin{frame}{Ejercicio 5: IS–LM con impuesto a la inversión}
\[
C = 1 + 0.8(1 - t)Y, \quad I = 2 - 0.4i, \quad G = 1, \quad t = 0.2
\]
\[
\frac{M}{P}=6, \quad L = 0.75Y - 1.5i
\]
\begin{enumerate}[label=\alph*)]
\item Encuentre \(Y\) e \(i\) de equilibrio.  
\item Si \(I_0\) cae 10\%, ¿cuánto debe subir \(G\)?  
\item Si \(I = 2 - 0.4(1+\phi)i\), con \(\phi=0.25\), obtenga el nuevo equilibrio.  
\item ¿Cómo debe variar \(M/P\) para volver a \(Y_a\)?  
\end{enumerate}
\espaciodesarrollo
\end{frame}

\begin{frame}{Ejercicio 5}
    
\end{frame}

% ===== Ejercicios extra =====
\begin{frame}{Ejercicio Extra 1}
Analice y grafique el efecto en el mercado del dinero de un aumento en la sensibilidad de la demanda de dinero al ingreso.  
Explique luego cómo se modifica el equilibrio IS–LM.
\espaciodesarrollo
\end{frame}

\begin{frame}{Ejercicio Extra 2}
\[
C = 420 + 0.7Y_D,\ t = 0.2,\ TR = 500,\ I = 4200 - 300i,\ G = 800,\ X_N = 400,
\]
\[
L = 10Y - 5000i,\quad \frac{M}{P} = 2000.
\]
\begin{itemize}
  \item Determine \(i\) y \(Y\) de equilibrio.  
  \item Evalúe el cambio en el déficit fiscal si las transferencias se duplican.  
\end{itemize}
\espaciodesarrollo
\end{frame}

% ===== Cierre =====
\begin{frame}{Síntesis}
\begin{itemize}
  \item El equilibrio IS–LM combina el mercado de bienes y el monetario.
  \item Las políticas fiscales y monetarias actúan desplazando las curvas.
  \item De este equilibrio surge la pendiente negativa de la DA.
  \item Comprender los desplazamientos permite anticipar efectos de políticas macroeconómicas.
\end{itemize}
\end{frame}

\end{document}
